% id: PHY-0034
% title: การกระจัดและระยะทางจากกราฟตำแหน่ง–เวลา
% camp: อื่น ๆ
% topic: การเคลื่อนที่แนวตรง
% year: 2549
% tags: สอวน., ม.4, การกระจัด, ระยะทาง, กราฟตำแหน่ง–เวลา
% origin: สอวน. ม.4 2549 · ข้อ 34 ในภาพที่ผู้ใช้ส่ง
% review: ถอดจากภาพแล้ว · รอตรวจทาน
% diagram: PHY-0034.svg
% scan: sources/PHY-0034-original.png
\documentclass[12pt]{article}
\usepackage{amsmath,amssymb,fontspec,polyglossia,tikz}
\setdefaultlanguage{thai}
\setmainfont{Noto Serif Thai}
\usetikzlibrary{arrows.meta}
\begin{document}
\textbf{ข้อ 34.} (สอวน. ม.4 2549)
อนุภาคหนึ่งกำลังเคลื่อนที่ในแนวเส้นตรง รูปข้างล่างแสดงให้เห็นกราฟตำแหน่งของอนุภาคเทียบกับจุดอ้างอิงที่เวลาต่าง ๆ โดยที่ได้เลือกเครื่องหมายบวกให้แทนทิศขวามือของจุดอ้างอิง

34.1 จงหาการกระจัดของอนุภาคนี้ตั้งแต่เวลา $t=0\ \mathrm{s}$ ไปจนถึงเวลา $t=7\ \mathrm{s}$

34.2 จงหาระยะทางที่อนุภาคเคลื่อนที่ในช่วงเวลานี้ด้วย

\begin{tikzpicture}[x=0.71571cm,y=0.55370cm,>=Stealth]
\draw[dashed,line width=.4pt] (1,0) -- (1,6.5);
\draw[dashed,line width=.4pt] (2,0) -- (2,6.5);
\draw[dashed,line width=.4pt] (3,0) -- (3,6.5);
\draw[dashed,line width=.4pt] (4,0) -- (4,6.5);
\draw[dashed,line width=.4pt] (5,0) -- (5,6.5);
\draw[dashed,line width=.4pt] (6,0) -- (6,6.5);
\draw[dashed,line width=.4pt] (7,0) -- (7,6.5);
\draw[dashed,line width=.4pt] (8,0) -- (8,6.5);
\draw[dashed,line width=.4pt] (0,1) -- (8.4,1);
\draw[dashed,line width=.4pt] (0,2) -- (8.4,2);
\draw[dashed,line width=.4pt] (0,3) -- (8.4,3);
\draw[dashed,line width=.4pt] (0,4) -- (8.4,4);
\draw[dashed,line width=.4pt] (0,5) -- (8.4,5);
\draw[dashed,line width=.4pt] (0,6) -- (8.4,6);
\draw[->,line width=.7pt] (-0.18,0) -- (8.7,0);
\draw[->,line width=.7pt] (0,-0.6) -- (0,6.6);
\draw[line width=.7pt] (1,-0.1204) -- (1,0.1204);
\node[anchor=center] at (1,-0.4816) {1};
\draw[line width=.7pt] (2,-0.1204) -- (2,0.1204);
\node[anchor=center] at (2,-0.4816) {2};
\draw[line width=.7pt] (3,-0.1204) -- (3,0.1204);
\node[anchor=center] at (3,-0.4816) {3};
\draw[line width=.7pt] (4,-0.1204) -- (4,0.1204);
\node[anchor=center] at (4,-0.4816) {4};
\draw[line width=.7pt] (5,-0.1204) -- (5,0.1204);
\node[anchor=center] at (5,-0.4816) {5};
\draw[line width=.7pt] (6,-0.1204) -- (6,0.1204);
\node[anchor=center] at (6,-0.4816) {6};
\draw[line width=.7pt] (7,-0.1204) -- (7,0.1204);
\node[anchor=center] at (7,-0.4816) {7};
\draw[line width=.7pt] (8,-0.1204) -- (8,0.1204);
\node[anchor=center] at (8,-0.4816) {8};
\draw[line width=.7pt] (-0.0931,1) -- (0,1);
\node[anchor=east] at (-0.2236,1) {1};
\draw[line width=.7pt] (-0.0931,2) -- (0,2);
\node[anchor=east] at (-0.2236,2) {2};
\draw[line width=.7pt] (-0.0931,3) -- (0,3);
\node[anchor=east] at (-0.2236,3) {3};
\draw[line width=.7pt] (-0.0931,4) -- (0,4);
\node[anchor=east] at (-0.2236,4) {4};
\draw[line width=.7pt] (-0.0931,5) -- (0,5);
\node[anchor=east] at (-0.2236,5) {5};
\draw[line width=.7pt] (-0.0931,6) -- (0,6);
\node[anchor=east] at (-0.2236,6) {6};
\node[anchor=east] at (-0.149,-0.3853) {0};
\node[anchor=west] at (8.849,-0.0482) {$t$ (s)};
\node[anchor=west] at (-0.149,6.9853) {$x$ (m)};
\draw[line width=.7pt] (0,1) -- (3,5) -- (4,2) -- (5,2) -- (7,3) -- (8,3);
\end{tikzpicture}
\end{document}
